\begin{frame}[t]{参考文献 5 / 5：本轮补充}
\fontsize{7.5}{10.5}\selectfont
\refentry{22}{Li et al. · Cell, 2026（本地正式版）}{UniPert-G2CP bridges genetic and chemical screens from molecular representation to phenotype modeling}{https://doi.org/10.1016/j.cell.2026.06.005}
\refentry{23}{Adduri et al. · Cell, 2026（本地正式版 / corrected proof）}{Predicting cellular responses to perturbation across diverse contexts with State}{https://doi.org/10.1016/j.cell.2026.07.052}
\refentry{24}{Sun et al. · Nature, 2026（本地正式版）}{An operational perturbation proteomics-based virtual cell model}{https://doi.org/10.1038/s41586-026-11001-9}
\refentry{25}{Burq et al. · bioRxiv, 2026-08-19（v1 预印本，未完成同行评审）}{Single-cell foundation models benefit from cross-modal training: adding proteomics data beats parameter scaling}{https://doi.org/10.64898/2026.08.14.744845}
\vspace{.18cm}
{\scriptsize\color{Muted}前三篇的数据规模以此次提供的 PDF 为准；第 25 篇包含正文与 SI，计数不一致及评估边界见 research/data-audit.md。}
\end{frame}
